% ---------------------------------------------------------------
% INTRODUCTION
% ---------------------------------------------------------------
\chapter*{Úvod}\addcontentsline{toc}{chapter}{Úvod}\markboth{Úvod}{Úvod}
\setcounter{page}{1}

Datová centra tvoří páteř moderního digitálního světa již několik desetiletí a jejich role se postupně vyvinula od jednoduchých úložišť dat až po komplexní výpočetní uzly v měřítku celých skladových hal \cite{barroso_data_2026}. Zlomovým okamžikem v jejich evoluci se stal nástup generativní umělé inteligence (AI), který po roce 2022 radikálně změnil nároky na infrastrukturu \cite{LBNL-2001637}. Na rozdíl od tradičních cloudových služeb vyžaduje trénování a provoz moderních AI modelů extrémní výpočetní výkon, což je spojeno s vysokými energetickými náklady \cite{bizo_silicon_2022}.

Současná AI infrastruktura čelí kritice kvůli svým environmentálním dopadům, které zahrnují nejen obrovskou spotřebu elektrické energie, ale také významné emise skleníkových plynů a narůstající odběr vody. Odhaduje se, že vývoj, trénování a provoz velkých jazykových modelů může během jednoho roku používání vyprodukovat až 160 000 tun uhlíku (jako provoz 30 000 benzinových aut) \cite{jegham_how_2025} a odebrat asi 5 miliard litrů sladké vody (což je polovina odběru Spojeného království) \cite{li_making_2025}. Tato technologická náročnost vytváří bezprecedentní tlak na lokální zdroje v místech, kde se AI datová centra koncentrují.

Ohniskem rozvoje AI datových center jsou v současnosti Spojené státy americké, kde se v nadcházejících letech očekává prudký nárůst poptávky po elektrické energii vyvolaný právě AI infrastrukturou \cite{chen_electricity_2025, davenport_ai_2024}. Tento trend se však postupně globalizuje a zasahuje i evropský region včetně České republiky. Vzhledem k tomu, že se otevření prvního specializovaného AI datového centra v České republice očekává na konci roku 2026 \cite{rabasova_u_2025} a největší datové centrum v ČR má být otevřeno na začátku roku 2027 \cite{cra_ceske_2025}, je nezbytné včas analyzovat rizika a příležitosti spojené s touto výstavbou.

Tato bakalářská práce slouží jako studie socioekonomických efektů investic do AI infrastruktury na území USA, které v této oblasti představují nejrozvinutější trh. Cílem práce je na základě analýzy amerických dat identifikovat klíčové korelace a vytvořit predikční model, který umožní simulovat dopady výstavby podobných center v České republice a potenciálně i v dalších zemích čelících AI transformaci. Cílem naopak není analyzovat výnosnost investic do umělé inteligence, nebo je hodnotit. Model má sloužit především k prozkoumání \textit{přímých} dopadů výstavby datových center, aniž by zahrnoval jejich produkty.